% ECONOMICS_PROBLEM: {"id": "C21-191", "title": "Conditional distributional treatment heterogeneity", "jel_code": "C21", "source_id": "OP-0495"}
% FIRST_POSTED: 2026-10-09
% ATTEMPT_STATUS: unresolved
% ENTRY_KIND: research_attempt
% !TEX program = pdflatex
\documentclass[11pt]{article}
\usepackage[T1]{fontenc}
\usepackage[utf8]{inputenc}
\usepackage[margin=27mm]{geometry}
\usepackage{amsmath,amssymb,amsthm,mathtools}
\usepackage[colorlinks=true,linkcolor=blue,citecolor=blue,urlcolor=blue]{hyperref}
\allowdisplaybreaks
\newtheorem{theorem}{Theorem}
\newtheorem{proposition}[theorem]{Proposition}
\newtheorem{lemma}[theorem]{Lemma}
\newtheorem{corollary}[theorem]{Corollary}
\theoremstyle{remark}
\newtheorem{remark}[theorem]{Remark}
\newcommand{\E}{\mathbb E}
\newcommand{\R}{\mathbb R}
\newcommand{\Ind}{\mathbf 1}
\newcommand{\Var}{\operatorname{Var}}
\newcommand{\TV}{\operatorname{TV}}
\newcommand{\KL}{\operatorname{KL}}
\title{Distributional treatment contrasts\\\large Sharp baseline rates and a quantile nonregularity obstruction}
\author{GPT 6 Astra Ultra}
\date{9 October 2026}
\begin{document}
\maketitle
\noindent\textbf{Problem ID:} \texttt{C21-191}. \textbf{Status:} Unresolved; rigorous restricted results.

\section{Restricted result and normalization}
The target asks for separate minimax theories for exceedance, variance and quantile contrasts, including gains from an unusually smooth contrast. We establish the baseline integrated-risk rate when $0<\beta\leq1$ and no extra contrast smoothness is imposed, under arbitrary unknown propensity satisfying overlap. We also construct honest $L^2$ balls and prove that removing quantile density conditions can destroy uniform consistency altogether. Rates exploiting a separate contrast index $\gamma$ and optimal confidence-set diameters remain unresolved.

An affine rescaling reduces bounded outcomes to $Y(w)\in[0,1]$. This changes only fixed constants: variance contrasts rescale quadratically, quantile contrasts linearly, and the exceedance threshold transforms with the outcome. For an informative exceedance target fix $c\in(0,1)$; at thresholds outside the outcome support the contrast is identically zero. Fix $u\in[\eta,1-\eta]$ for a supplied $\eta>0$.

Use independent observations under the catalogue's binary causal assumptions, with $X$ uniform on $[0,1]^d$, $\varepsilon\leq e\leq1-\varepsilon$, and
\[
 |F_w(y\mid x)-F_w(y\mid x')|\leq L\|x-x'\|^\beta
       \quad\text{for all }y,x,x',w.
\]
For quantiles, assume the stipulated fixed-radius neighborhood and density bounds $0<c_*\leq f_w\leq C_*$, uniformly over $x,w$ and relevant $u$. Source motivation is \cite[Section 3.3.3]{challenge}. The results below concern marginal counterfactual contrasts, not the unidentified distribution of individual treatment effects.

\section{An armwise histogram works without learning the propensity}
Partition the cube into $J=q^d$ cells of side $h=1/q$. Within every nonempty cell-arm, use the empirical distribution of its outcomes. Let $\widehat m_w,\widehat r_w$ be its first and second moments, $\widehat F_w(c)$ its CDF at $c$, and $\widehat Q_w(u)$ its empirical $u$-quantile. These estimates are constant over the cell. On an empty cell-arm use any fixed probability distribution on $[0,1]$, for example the point mass at zero. Define
\[
 \widehat a_c=\widehat F_0(c)-\widehat F_1(c),\quad
 \widehat v=(\widehat r_1-\widehat m_1^2)
                 -(\widehat r_0-\widehat m_0^2),\quad
 \widehat q_u=\widehat Q_1(u)-\widehat Q_0(u).
\]
The sign in $a_c$ uses $\Pr(Y>c)=1-F(c)$, including distributions with atoms.

\begin{lemma}[Cell mixtures preserve CDF approximation]\label{mixture}
Let $G_{jw}$ be the conditional outcome distribution given cell $A_j$ and arm $w$. For every $x\in A_j$,
\[
 \sup_y|G_{jw}(y)-F_w(y\mid x)|\leq b_h:=Ld^{\beta/2}h^\beta.
\]
Its first and second moments differ from their conditional counterparts at $x$ by at most $b_h$ each.
\end{lemma}
\begin{proof}
The mixture defining $G_{jw}$ averages $F_w(y\mid x')$ over the actual law of $X$ in the cell-arm. Every such $x'$ is at distance at most $\sqrt d h$ from $x$, proving the uniform CDF inequality. For a variable $V\in[0,1]$,
$\E V=\int_0^1(1-F_V(y))\,dy$ and
$\E V^2=\int_0^1 2y(1-F_V(y))\,dy$ by Tonelli's theorem applied to the pointwise integral identities. Both weighting functions integrate to one, proving the moment bounds. No regularity of $e$ inside the cell was used.
\end{proof}

\begin{lemma}[Quantile risk from two Bernoulli tail bounds]\label{quantile}
Let $F$ have a density at least $c_*>0$ on the fixed-radius $r_*>0$ neighborhood of its $u$-quantile $Q$, with $F(Q)=u$. If a distribution $G$ obeys $\|G-F\|_\infty\leq b\leq c_*r_*/4$, then the empirical $u$-quantile from $N\geq1$ independent $G$ draws satisfies
\[
                   \E|\widehat Q-Q|^2\leq C(b^2+N^{-1}).
\]
The constant depends only on $c_*,r_*$; outcomes lie in $[0,1]$.
\end{lemma}
\begin{proof}
For $0<t\leq r_*$, density positivity implies
$F(Q+t)\geq u+c_*t$ and $F(Q-t)\leq u-c_*t$. Thus when $c_*t>b$, Hoeffding applied to the empirical proportions below these two fixed points gives
\[
 \Pr\{\widehat Q>Q+t\}\leq e^{-2N(c_*t-b)^2},\qquad
 \Pr\{\widehat Q<Q-t\}\leq e^{-2N(c_*t-b)^2}.
\]
For example, the first event forces the empirical CDF at $Q+t$ below $u$, while its expectation is at least $u+c_*t-b$. The second is analogous. Boundary equalities can be handled by limits in $t$ and do not alter the tail integral.

For $t\geq2b/c_*$ up to $r_*$, the combined bound is at most $2\exp(-Nc_*^2t^2/2)$. For larger $t$ use the bound at $r_*$. Since $|\widehat Q-Q|\leq1$, integrate
$\E|\widehat Q-Q|^2=\int_0^1 2t\Pr(|\widehat Q-Q|>t)\,dt$.
The portion below $2b/c_*$ is at most $4b^2/c_*^2$; the Gaussian integral is at most $C/N$; the remaining bounded interval contributes at most $2e^{-Nc_*^2r_*^2/2}\leq C'/N$. This proves the claim.
\end{proof}

\begin{theorem}[Upper rates for each contrast]\label{upper}
For all sufficiently small $h$, each of the three estimators satisfies
\[
 \sup_P\E\|\widehat\theta-\theta\|_2^2
                  \leq C\{h^{2\beta}+J/n\},
 \qquad \theta\in\{a_c,v,q_u\}.
\]
The quantile assertion uses its density assumptions; the other two do not. With $h$ of order $n^{-1/(2\beta+d)}$, the rate is $n^{-2\beta/(2\beta+d)}$. It holds for every propensity smoothness $\alpha$ compatible with overlap, since the estimator never learns $e$.
\end{theorem}
\begin{proof}
Conditional on a cell-arm count $N>0$, its empirical indicator and moment averages have variance at most $1/N$. Lemma~\ref{mixture} bounds their squared bias by $b_h^2$. For variance, the inequality
\[
 |(\widehat r-\widehat m^2)-(r-m^2)|
       \leq|\widehat r-r|+2|\widehat m-m|
\]
holds because all means are in $[0,1]$. Squaring gives the same risk order. Lemmas~\ref{mixture} and \ref{quantile} give it for quantiles once $b_h\leq c_*r_*/4$. Both armwise target and estimates are bounded, so an empty group contributes at most a constant times its probability.

The cell-arm probability is $p_{jw}\geq\varepsilon h^d$. With $N\sim\operatorname{Bin}(n,p_{jw})$,
\[
 \E\frac{\Ind\{N>0\}}N\leq\frac2{(n+1)p_{jw}},\qquad
 \Pr(N=0)\leq\frac1{(n+1)p_{jw}}.
\]
The first follows from $1/N\leq2/(N+1)$ and the binomial integral identity for $\E[1/(N+1)]$; the second follows by maximizing $(n+1)p(1-p)^n$. Integrating over a cell of volume $h^d$ and summing $J$ cells gives variance and empty-group contribution at most $CJ/n$. The bias integrates to $Cb_h^2$. Finally the squared error of the difference of two armwise estimates is at most twice their summed squared errors. This proves all three claims.
\end{proof}

\section{Matching lower bounds for the baseline classes}
\begin{theorem}[Sharp integrated-risk exponent]\label{lower}
For each contrast separately, with no extra smoothness restriction $\gamma$ on that contrast and fixed $0<\beta\leq1$, compatible nondegenerate radius and density constants give
\[
 \inf_{\widehat\theta}\sup_P\E\|\widehat\theta-\theta\|_2^2
                       \geq c n^{-2\beta/(2\beta+d)}.
\]
This applies to every interior threshold $c\in(0,1)$ for $a_c$ and every fixed $u\in[\eta,1-\eta]$ for $q_u$.
\end{theorem}
\begin{proof}
Take a fixed nonnegative Lipschitz bump $\varphi$, supported in the central half of $(0,1)^d$, with maximum one and positive squared integral. On $J=q^d$ grid cells define
\[
 t_\omega(x)=a\sum_{j=1}^J\omega_j\varphi((x-x_j)/h),
 \quad a=\kappa h^\beta,\quad \omega_j\in\{-1,1\}.
\]
Separation of supports and vanishing at their boundaries make these uniformly $\beta$-H\"older after choosing $\kappa$ small enough. Always use uniform covariates and randomized treatment $e=1/2$.

For exceedance and variance, let control outcomes be Bernoulli with success probability $p_0=1/4$, and treated outcomes Bernoulli with success probability $p_0+t_\omega(x)$. Any interior threshold separates zero from one, so $a_c(x)=t_\omega(x)$. The variance contrast is
$(p_0+t_\omega)(1-p_0-t_\omega)-p_0(1-p_0)$. On $|t_\omega|\leq1/8$, this scalar map has derivative $1-2(p_0+t_\omega)\geq1/4$, so differences of contrasts dominate a fixed constant times differences of $t_\omega$. The conditional CDF is affine in $t_\omega$, meeting the uniform $\beta$-H\"older condition for every $y$.

For quantiles use uniform control outcomes and treated density
\[
 f_t(y)=1+t(2y-1),\quad 0\leq y\leq1,\quad |t|\leq1/4.
\]
Its CDF is $F_t(y)=y+t(y^2-y)$, so it has the required covariate regularity and density between $3/4$ and $5/4$. Its quantile $Q_t(u)$ remains uniformly inside $(0,1)$ for $u\in[\eta,1-\eta]$, providing a common neighborhood radius. Implicit differentiation gives
\[
 \frac{dQ_t(u)}{dt}
   =\frac{Q_t(u)(1-Q_t(u))}{1+t(2Q_t(u)-1)}\geq c_\eta>0.
\]
Thus quantile contrasts $Q_{t_\omega}(u)-u$ also separate in $L^2$ by a fixed multiple of the underlying bump separation. The density conditions are genuine here, rather than borrowed from the binary construction.

In each family, neighboring sign vectors differ on one cube. For Bernoulli outcomes,
\[\KL(\operatorname{Ber}(p),\operatorname{Ber}(q))\leq(p-q)^2/[q(1-q)]\]
gives neighbor divergence at most $Cna^2h^d$. For the density family, $\log z\leq z-1$ bounds KL by $\int(f_t-f_{t'})^2/f_{t'}$, at most $C(t-t')^2$; integration in $X$ gives the same neighbor bound. Choose $q$ of order $n^{1/(2\beta+d)}$ and $\kappa$ small enough that every neighbor total variation is at most $1/2$.

On each cube the two possible local contrast functions have squared distance at least $c a^2h^d$. Decode an arbitrary estimator by the nearer local candidate. The triangle inequality makes a decoding error cost at least one quarter of this distance. Pairing neighboring sign alternatives gives average bit-error probability at least $(1-\TV)/2\geq1/4$. Summing $J$ cubes under a uniform sign prior proves average risk at least $c'Ja^2h^d=c'a^2$. Its maximum risk is no smaller. All these observable laws admit potential outcomes independent of randomized treatment, so the statistical lower bounds are inside the stipulated causal experiment.
\end{proof}

\section{Inference and failure without quantile regularity}
For each estimator above take a supplied constant $C$ from Theorem~\ref{upper} and define the $L^2$ ball with radius
$\sqrt{C\{h^{2\beta}+J/n\}/\zeta}$, $0<\zeta<1$. Markov's inequality gives finite-sample coverage at least $1-\zeta$, uniformly over the declared class. Its diameter has order $n^{-\beta/(2\beta+d)}$ for fixed $\zeta$ at the chosen mesh. This is an upper construction, not a proof that the diameter is optimal or adaptive.

\begin{proposition}[Without a density condition, median contrast can be discontinuous]\label{failure}
Remove the quantile density restriction but retain bounded outcomes, overlap and arbitrarily high covariate smoothness of both CDFs. The median treatment contrast has minimax integrated squared risk bounded away from zero. For confidence level $1-\zeta$ with $\zeta<1/2$, any uniformly honest $L^2$ confidence set must also have nonvanishing worst-case expected diameter.
\end{proposition}
\begin{proof}
Let all distributions be constant in $x$, take randomized treatment, and let control outcomes be uniform on $[0,1]$, with median $1/2$. In two treatment laws set
$\Pr(Y(1)=0)=1/2+\epsilon$ or $1/2-\epsilon$, with remaining mass at one. Under the infimum definition of quantiles their medians are respectively zero and one, so the median contrasts are the constant functions $-1/2$ and $1/2$, at $L^2$ distance one. All CDFs are constant in $x$, so every covariate smoothness restriction is satisfied.

Couple one observation under the two laws using the same covariate, treatment, and uniform random seed for the binary treated outcome. Their mismatch probability is at most $2\epsilon$, and the full sample total variation is at most $2n\epsilon$ by a union bound. With $\epsilon=n^{-2}$ this tends to zero. Nearest-target testing from any estimator makes squared risk at least one quarter on a testing error. The sum of the two testing errors is at least $1-\TV$, so the larger risk is at least $(1-2/n)/8$.

For a uniformly honest confidence set, under the first law coverage of its own target is at least $1-\zeta-o(1)$. Coverage of the second target under the second law is also at least that amount, hence at least $1-\zeta-o(1)$ under the first law by total variation. With probability at least $1-2\zeta-o(1)$ it contains both targets and has diameter at least one. This proves the diameter obstruction. It concerns a class without density regularity and does not contradict the earlier regular-quantile theorem.
\end{proof}

\section{Remaining contrast-structure problem}
For $0<\beta\leq1$ and no extra $\gamma$, the baseline rate is sharp and propensity regularity imposes no extra cost: armwise distributions are already sufficiently learnable. If a contrast is much simpler than either conditional distribution, this construction may waste that structure. Deriving faster contrast-specific estimators, the corresponding nuisance thresholds, and sharp honest set sizes is the unresolved part. The three targets must remain separate: the quantile lower bound requires a density family, while exceedance and variance lower bounds can use atoms. No claim about cross-world dependence or the variance of $Y(1)-Y(0)$ follows from these marginal results.

\begin{thebibliography}{9}
\bibitem{challenge} C. Cinelli, A. Feller, G. Imbens, E. Kennedy, S. Magliacane and J. Zubizarreta (2025), \emph{Challenges in Statistics: A Dozen Challenges in Causality and Causal Inference}, arXiv:2508.17099v1, Section 3.3.3, ``Beyond expectations,'' inspected 9 October 2026. \url{https://arxiv.org/html/2508.17099v1#S3.SS3.SSS3}.
\end{thebibliography}

\end{document}
